%%% -*- coding: utf-8 -*-
% !TeX root = main-slides.tex

\lecture{Overview}{overview}
\section{Overview}

\showtitlepage{Goals and Content}

\begin{frame}{Goals}
  The goals of this tutorial is to provide the necessary skills for \dots
  \begin{itemize}
    \item programming C\# to the extent required for Unity, assuming
      programming skills in Java,
    \item developing scripts for Unity in general,
    \item enhancing SEE\footnote{\url{https://github.com/uni-bremen-agst/SEE}} in particular.
  \end{itemize}
  \vspace{3ex}

  Things not covered:
  \begin{itemize}
    \item all the latest features of C\#\footnote{Unity 6 is using C\# version 9.0},
    \item all features of Unity, e.g., animation, textures, shaders, etc.
    \item complete insights of SEE's implementation,
    \item all third-party components SEE is using.
  \end{itemize}
\end{frame}

\begin{frame}{How This Course Is Graded}
  There is no written exam. You work on a small programming task and then
  explain what you did.
  \begin{itemize}
    \item \textbf{Friday afternoon, two and a half hours, in this room.} You
      receive a small program built on SEE's graph model with several planted
      defects. Find them, write a failing test for each, repair them. Your
      defects are yours alone.
    \item \textbf{Handing in is a push.} At the end of the session you push
      your branch and we note the commit. Nothing changes after that.
    \item \textbf{Defense in the week after, fifteen minutes each.} You walk us
      through your work, answer questions on it, and make one small change
      while we watch.
    \item \textbf{Forty percent the code, sixty percent the defense.} Finishing
      everything on Friday matters less than being able to account for what you
      did.
    \item \textbf{AI assistants are allowed.} Say in your commit message what
      you used them for. You will be asked about the result either way.
  \end{itemize}
  \vspace{1ex}

  On Thursday we build the assignment repository together, so that Friday
  starts with a working machine.
\end{frame}

%%% Local Variables:
%%% mode: latex
%%% TeX-master: "main-slides"
%%% mode: reftex
%%% mode: flyspell
%%% End:
